\documentclass[../../main.tex]{subfiles}% 注意这里的文件路径不能用 ./main.tex ,否则用latexmk编译子文件会报错
\graphicspath{{\subfix{./image/}}{\subfix{../../image/}}} % 指定图片引用路径,后续可以直接使用图片文件名
% 如果图片存放在上级目录,则使用 ../../image/ 作为备用路径

% 例如:
% \begin{figure}[H]
% \centering
% \includegraphics[width=0.3\textwidth]{图.png}
% \caption{}
% \label{figure:图}
% \end{figure}
% 注意:上述\label{}一定要放在\caption{}之后,否则引用图片序号会只会显示??.

\begin{document}

\section{极值原理和最大模估计}

\subsection{极值原理}

对于一个物体，如果其内部没有热源，则在整个热传导过程中，温度总是趋于平衡：温度高处的热量向温度低处传递，从而导致温度高处的热量减少，温度降低；而温度低处从温度高处吸收热量，从而导致温度低处的热量增加，温度升高。正是因为这样的物理效应，物体在一段时间内的最高温度和最低温度不可能在物体的内部达到，因而只能在初始时刻或物体的边界上达到。这种物理现象我们在数学上称为\textbf{极值原理}。

\begin{definition}
记区域 $Q_T=(0,l)\times(0,T]$,其中$l,T>0$.$Q_T$ 的侧边和底边我们称为\textbf{抛物边界}，通常记为 $\partial_p Q_T$。实际上
\begin{align*}
\partial_p Q_T=\partial Q_T\setminus(0,l)\times\{T\}.
\end{align*}
\end{definition}

\begin{theorem}[极值原理]\label{theorem:PDE极值原理}
假设 $u\in C^{2,1}(Q_T)\cap C(\overline{Q}_T)$ 满足方程 \begin{align*}
\mathcal{L}u=u_t-a^2u_{xx}=f(x,t)\leqslant 0,
\end{align*}
则 $u(x,t)$ 在 $\overline{Q}_T$ 上的最大值必在抛物边界 $\partial_p Q_T$ 上达到，即
\begin{align*}
\max_{\overline{Q}_T}u(x,t)=\max_{\partial_p Q_T}u(x,t).
\end{align*}
\end{theorem}
\begin{proof}
由于 $\partial_p Q_T\subset\overline{Q}_T$，因此不等式
\begin{align*}
\max_{\overline{Q}_T}u(x,t)\geqslant \max_{\partial_p Q_T}u(x,t)
\end{align*}
成立。我们只要证明不等式
\begin{align*}
\max_{\overline{Q}_T}u(x,t)\leqslant \max_{\partial_p Q_T}u(x,t)
\end{align*}
就完成了定理的证明。

(1) 当 $f<0$ 时，我们断言 $u$ 在 $\overline{Q}_T$ 上的最大值不能在 $Q_T$ 内达到。否则，存在一点 $(x_0,t_0)\in Q_T$ 使得
\begin{align*}
u(x_0,t_0)=\max_{\overline{Q}_T}u(x,t).
\end{align*}
由微积分的定理我们得到
\begin{align*}
u_x(x_0,t_0)=0,\qquad u_{xx}(x_0,t_0)\leqslant 0,
\end{align*}
且
\begin{align*}
u_t(x_0,t_0)=0,\quad t_0<T,\\
u_t(x_0,t_0)\geqslant 0,\quad t_0=T.
\end{align*}
因此
\begin{align*}
f(x_0,t_0)=u_t(x_0,t_0)-a^2u_{xx}(x_0,t_0)\geqslant 0.
\end{align*}
这与假设 $f<0$ 矛盾，从而 $u$ 不可能在 $Q_T$ 内达到 $\overline{Q}_T$ 上的最大值，因而只能在抛物边界 $\partial_p Q_T$ 上达到最大值。于是所要证明的不等式成立。

(2) 当 $f\leqslant 0$ 时，我们构造辅助函数，将证明归结到上面的情形。为此，对于任意 $\varepsilon>0$，考虑辅助函数
\begin{align*}
v(x,t)=u(x,t)-\varepsilon t.
\end{align*}
计算得到
\begin{align*}
\mathcal{L}v=\mathcal{L}u-\varepsilon=f-\varepsilon<0.
\end{align*}
由 (1) 的断言，$v(x,t)$ 在 $\overline{Q}_T$ 上的最大值不可能在 $Q_T$ 内达到，因此
\begin{align*}
\max_{\overline{Q}_T}v(x,t)=\max_{\partial_p Q_T}v(x,t).
\end{align*}
于是
\begin{align*}
\max_{\overline{Q}_T}u(x,t)\leqslant \max_{\overline{Q}_T}v(x,t)+\varepsilon T 
\leqslant \max_{\partial_p Q_T}v(x,t)+\varepsilon T 
\leqslant \max_{\partial_p Q_T}u(x,t)+\varepsilon T.
\end{align*}
令 $\varepsilon\to 0$，则得所要证明的不等式。至此完成定理的证明.
\end{proof}

\begin{corollary}
\begin{enumerate}[(1)]
\item 假设 $u\in C^{2,1}(Q_T)\cap C(\overline{Q}_T)$ 满足方程
\begin{align*}
\mathcal{L}u=u_t-a^2u_{xx}=f(x,t)\leqslant 0,
\end{align*}
则 $u(x,t)$ 在 $\overline{Q}_T$ 上的最小值必在抛物边界 $\partial_p Q_T$ 达到，即
\begin{align*}
\min_{\overline{Q}_T}u(x,t)=\min_{\partial_p Q_T}u(x,t);
\end{align*}

\item 假设 $u\in C^{2,1}(Q_T)\cap C(\overline{Q}_T)$ 满足方程
\begin{align*}
\mathcal{L}u=u_t-a^2u_{xx}=f(x,t) = 0,
\end{align*}
则 $u(x,t)$ 在 $\overline{Q}_T$ 上的最大值和最小值必在抛物边界 $\partial_p Q_T$ 达到。
\end{enumerate}
\end{corollary}
\begin{proof}
考虑 $v(x,t)=-u(x,t)$ 满足的方程，应用\hyperref[theorem:PDE极值原理]{极值原理}即得证.
\end{proof}

\begin{corollary}[比较原理]\label{corollary:PDE比较原理}
假设 $u,v\in C^{2,1}(Q_T)\cap C(\overline{Q}_T)$ 满足 $\mathcal{L}u\leqslant \mathcal{L}v$，$u|_{\partial_p Q_T}\leqslant v|_{\partial_p Q_T}$，则在 $\overline{Q}_T$ 上 $u(x,t)\leqslant v(x,t)$。
\end{corollary}
\begin{proof}
考虑 $w(x,t)=u(x,t)-v(x,t)$。由于 $\mathcal{L}w\leqslant 0$，对 $w(x,t)$ 应用\hyperref[theorem:PDE极值原理]{极值原理}得到
\begin{align*}
\max_{\overline{Q}_T}w(x,t)\leqslant \max_{\partial_p Q_T}w(x,t)\leqslant 0.
\end{align*}
至此我们完成比较原理的证明.
\end{proof}

\subsection{第一边值问题的最大模估计}

\begin{definition}[热方程第一边值问题]
\textbf{热方程第一边值问题}定义为
\begin{align}
\begin{cases}
\mathcal{L}u=u_t-a^2u_{xx}=f(x,t), & (x,t)\in Q_T,\\
u|_{t=0}=\varphi(x), & x\in[0,l],\\
u|_{x=0}=g_1(t),\quad u|_{x=l}=g_2(t), & t\in[0,T].
\end{cases} \label{eq:heat-dirichlet-max}
\end{align}
其中$u=u(x,t)$ 为未知函数，$a>0$ 为常数，$f(x,t)$,$\varphi(x)$,$g_1(t),g_2(t)$ 均为已知函数,且 $Q_T=(0,l)\times(0,T]$。
\end{definition}

\begin{theorem}\label{theorem:热方程第一边值问题的最大模估计}
假设 $u\in C^{2,1}(Q_T)\cap C(\overline{Q}_T)$ 是热方程第一边值问题\eqref{eq:heat-dirichlet-max} 的解，则
\begin{align}
\max_{\overline{Q}_T}|u(x,t)|\leqslant FT+B, \label{eq:heat-max-est1}
\end{align}
其中
\begin{align*}
F=\sup_{Q_T}|f|,\qquad B=\max\left\{\max_{[0,l]}|\varphi|,\ \max_{[0,T]}|g_1|,\ \max_{[0,T]}|g_2|\right\}.
\end{align*}
\end{theorem}
\begin{proof}
考虑辅助函数
\begin{align*}
w(x,t)=Ft+B\pm u(x,t).
\end{align*}
容易验证
\begin{align*}
\mathcal{L}w=F\pm f\geqslant 0,\qquad w|_{\partial_p Q_T}\geqslant B\pm u\geqslant 0.
\end{align*}
由\hyperref[theorem:PDE极值原理]{极值原理}，在 $\overline{Q}_T$ 上 $w(x,t)\geqslant 0$，从而
\begin{align*}
|u(x,t)|\leqslant FT+B,\quad (x,t)\in\overline{Q}_T.
\end{align*}
两端取上确界，则定理得证.
\end{proof}

\begin{corollary}
热方程第一边值问题\eqref{eq:heat-dirichlet-max} 的解在 $C^{2,1}(Q_T)\cap C(\overline{Q}_T)$ 中是惟一的。
\end{corollary}
\begin{proof}
由于热方程第一边值问题\eqref{eq:heat-dirichlet-max} 是线性的，为证明惟一性，只需证明当 $f\equiv 0,\varphi\equiv 0,g_1=g_2\equiv 0$ 时，问题只有零解。由\hyperref[theorem:热方程第一边值问题的最大模估计]{最大模估计}，这是显然的.
\end{proof}

\begin{corollary}
热方程第一边值问题\eqref{eq:heat-dirichlet-max} 在 $C^{2,1}(Q_T)\cap C(\overline{Q}_T)$ 中的解连续地依赖于非齐次项 $f$，初值 $\varphi$，边值 $g_1,g_2$。
\end{corollary}
\begin{remark}
由上述推论我们知道最大模估计蕴涵着解的惟一性和稳定性。
\end{remark}


\subsection{第二、第三边值问题的最大模估计}

\begin{lemma}\label{lemma:热方程的混合边值问题引理}
假设 $u\in C^{2,1}(Q_T)\cap C^{1,0}(\overline{Q}_T)$ 是热方程的混合边值问题\eqref{eq:heat-third-max}
\begin{align}
\begin{cases}
\mathcal{L}u=u_t-a^2u_{xx}=f(x,t), & (x,t)\in Q_T,\\
u|_{t=0}=\varphi(x), & x\in[0,l],\\
[-u_x+\alpha(t)u]|_{x=0}=g_1(t), & t\in[0,T],\\
[u_x+\beta(t)u]|_{x=l}=g_2(t), & t\in[0,T],
\end{cases} \label{eq:heat-third-max}
\end{align}
的解,
则在 $\overline{Q}_T$ 上 $u(x,t)\geqslant 0$。
\end{lemma}
\begin{proof}

(1) 先假设 $u(x,t)$ 满足边值条件
\begin{align*}
[-u_x+\alpha(t)u]|_{x=0}>0,\quad t\in[0,T],\\
[u_x+\beta(t)u]|_{x=l}>0,\quad t\in[0,T].
\end{align*}
由极值原理我们知道，$u(x,t)$ 在 $\overline{Q}_T$ 的最小值在抛物边界 $\partial_p Q_T$ 上达到。此时我们只需证明 $u(x,t)$ 在抛物边界 $\partial_p Q_T$ 上的最小值一定非负，从而在 $\overline{Q}_T$ 上 $u(x,t)\geqslant 0$。

如果 $u(x,t)$ 在初始时刻 $t=0$ 达到最小值，显然最小值非负。我们只需说明 $u$ 不可能在边界 $x=0$ 和 $x=l$ 上达到负的最小值。如果 $u(x,t)$ 在某点 $(0,t_0)$ 达到负的最小值，则
\begin{align*}
-u_x(0,t_0)\leqslant 0,\qquad \alpha(t_0)u(0,t_0)\leqslant 0,
\end{align*}
这与假设矛盾。同理 $u(x,t)$ 不可能在边界 $x=l$ 上达到负的最小值。这就说明，如果 $u$ 在边界 $x=0$ 和 $x=l$ 上达到最小值，则最小值一定非负。因此 $u(x,t)$ 在抛物边界 $\partial_p Q_T$ 上的最小值一定非负，从而在 $\overline{Q}_T$ 上 $u(x,t)\geqslant 0$。

(2) 对于一般情形，我们将构造辅助函数，把证明归结到上面的情形。为此，对于任意 $\varepsilon>0$，我们考虑辅助函数
\begin{align*}
v(x,t)=u(x,t)+\varepsilon w(x,t),
\end{align*}
其中
\begin{align}
w(x,t)=2a^2t+\left(x-\frac{l}{2}\right)^2. \label{eq:aux-w}
\end{align}
注意到多项式函数 $w$ 满足 $\mathcal{L}w=0$，我们计算得知
\begin{align*}
\mathcal{L}v=\mathcal{L}u\geqslant 0,\quad (x,t)\in Q_T.
\end{align*}
显然
\begin{align*}
v|_{t=0}=u|_{t=0}+\varepsilon\left(x-\frac{l}{2}\right)^2\geqslant 0,\quad x\in[0,l].
\end{align*}
由混合问题的边值条件我们得到，当 $t\in[0,T]$ 时，
\begin{align*}
[-v_x+\alpha(t)v]|_{x=0}
&=[-u_x+\alpha(t)u]|_{x=0}
+\varepsilon\left[l+\alpha(t)\left(2a^2t+\frac{l^2}{4}\right)\right]>0,\\
[v_x+\beta(t)v]|_{x=l}
&=[u_x+\beta(t)u]|_{x=l}
+\varepsilon\left[l+\beta(t)\left(2a^2t+\frac{l^2}{4}\right)\right]>0.
\end{align*}
应用 (1) 已证的结论我们得到，在 $\overline{Q}_T$ 上 $v(x,t)\geqslant 0$，即
\begin{align*}
u(x,t)\geqslant -\varepsilon\left[2a^2t+\left(x-\frac{l}{2}\right)^2\right],\quad (x,t)\in\overline{Q}_T.
\end{align*}
令 $\varepsilon\to 0$，则在 $\overline{Q}_T$ 上 $u(x,t)\geqslant 0$。这就是要证明的结论.
\end{proof}

\begin{theorem}
假设 $u\in C^{2,1}(Q_T)\cap C^{1,0}(\overline{Q}_T)$ 是热方程的混合边值问题\eqref{eq:heat-third-max1}
\begin{align}
\begin{cases}
\mathcal{L}u=u_t-a^2u_{xx}=f(x,t), & (x,t)\in Q_T,\\
u|_{t=0}=\varphi(x), & x\in[0,l],\\
[-u_x+\alpha(t)u]|_{x=0}=g_1(t), & t\in[0,T],\\
[u_x+\beta(t)u]|_{x=l}=g_2(t), & t\in[0,T],
\end{cases} \label{eq:heat-third-max1}
\end{align}
的解，则
\begin{align}
\max_{\overline{Q}_T}|u(x,t)|\leqslant C(F+B), \label{eq:heat-max-est2}
\end{align}
其中常数 $C$ 只依赖于 $a,l$ 和 $T$，而
\begin{align*}
F=\sup_{Q_T}|f|,\qquad B=\max\left\{\max_{[0,l]}|\varphi|,\ \max_{[0,T]}|g_1|,\ \max_{[0,T]}|g_2|\right\}.
\end{align*}
\end{theorem}
\begin{proof}
我们构造辅助函数
\begin{align*}
v(x,t)=Ft+Bz(x,t)\pm u(x,t),
\end{align*}
其中 $F,B$ 在定理中定义，且
\begin{align*}
z(x,t)=1+\frac{1}{l}w(x,t),
\end{align*}
而 $w(x,t)$ 由 \eqref{eq:aux-w} 定义。不难验证
\begin{align*}
\begin{cases}
\mathcal{L}v=F\pm f(x,t)\geqslant 0, & (x,t)\in Q_T,\\
v|_{t=0}\geqslant B\pm\varphi(x)\geqslant 0, & x\in[0,l],\\
[-v_x+\alpha(t)v]|_{x=0}\geqslant B\pm g_1(t)\geqslant 0, & t\in[0,T],\\
[v_x+\beta(t)v]|_{x=l}\geqslant B\pm g_2(t)\geqslant 0, & t\in[0,T].
\end{cases}
\end{align*}
我们利用\cref{lemma:热方程的混合边值问题引理}得到，在 $\overline{Q}_T$ 上 $v(x,t)\geqslant 0$。于是，对于 $(x,t)\in\overline{Q}_T$，有
\begin{align*}
|u(x,t)|\leqslant Ft+Bz(x,t) 
\leqslant FT+\left(1+\frac{2a^2T}{l}+\frac{l}{4}\right)B.
\end{align*}
令
\begin{align*}
C=\max\left\{T,\ 1+\frac{2a^2T}{l}+\frac{l}{4}\right\},
\end{align*}
上式两端取上确界我们立即得到估计 \eqref{eq:heat-max-est2}.
\end{proof}


\subsection{初值问题的最大模估计}

\begin{theorem}
假设 $u\in C^{2,1}(Q_T)\cap C(\overline{Q}_T)$ 是初值问题 \eqref{eq:heat-cauchy-max}
\begin{align}
\begin{cases}
\mathcal{L}u=u_t-a^2u_{xx}=f(x,t), & (x,t)\in Q_T,\\
u(x,0)=\varphi(x), & x\in\mathbb{R}.
\end{cases} \label{eq:heat-cauchy-max}
\end{align}
的有界解，则
\begin{align}
\sup_{\overline{Q}_T}|u(x,t)|\leqslant T\sup_{Q_T}|f(x,t)|+\sup_{\mathbb{R}}|\varphi(x)|. \label{eq:heat-max-cauchy}
\end{align}
\end{theorem}
\begin{remark}
定理中关于 $u(x,t)$ 有界的假设可以改进为
\begin{align}
|u(x,t)|\leqslant Me^{Ax^2},\quad (x,t)\in Q_T, \label{eq:growth}
\end{align}
这里 $M,A$ 是正常数。通常称条件 \eqref{eq:growth} 为增长性条件。
\end{remark}
\begin{proof}
对于任意 $L>0$，考虑区域 $Q_T^L=(-L,L)\times(0,T]$，并记
\begin{align*}
F=\sup_{Q_T}|f(x,t)|,\qquad \Phi=\sup_{\mathbb{R}}|\varphi(x)|.
\end{align*}
令
\begin{align*}
M=\sup_{\overline{Q}_T}|u(x,t)|.
\end{align*}
在 $Q_T^L$ 上考虑辅助函数
\begin{align*}
w(x,t)=Ft+\Phi+v_L(x,t)\pm u(x,t),
\end{align*}
其中
\begin{align*}
v_L(x,t)=\frac{M}{L^2}(x^2+2a^2t).
\end{align*}
我们计算得到
\begin{align*}
\begin{cases}
\mathcal{L}w=F\pm f(x,t)\geqslant 0, & (x,t)\in Q_T^L,\\
w|_{t=0}\geqslant \Phi\pm\varphi(x)\geqslant 0, & x\in[-L,L],\\
w|_{x=\pm L}\geqslant M\pm u\geqslant 0, & t\in[0,T].
\end{cases}
\end{align*}
在 $Q_T^L$ 上利用\hyperref[theorem:PDE极值原理]{极值原理}可得 $\min_{\overline{Q}_T^L}w(x,t)\geqslant 0$。而对于任意 $(x_0,t_0)\in Q_T$，存在足够大的 $L$，使得 $(x_0,t_0)\in Q_T^L$。由 $w(x_0,t_0)\geqslant 0$ 得到
\begin{align*}
|u(x_0,t_0)|\leqslant Ft_0+\Phi+\frac{M}{L^2}(x_0^2+2a^2t_0).
\end{align*}
令 $L\to+\infty$，则
\begin{align*}
|u(x_0,t_0)|\leqslant Ft_0+\Phi\leqslant FT+\Phi.
\end{align*}
由 $(x_0,t_0)$ 的任意性，我们完成定理的证明.
\end{proof}

\begin{theorem}
假设 $u\in C^{2,1}(Q_T)\cap C(\overline{Q}_T)$ 是初值问题 \eqref{eq:heat-cauchy-max1}
\begin{align}
\begin{cases}
\mathcal{L}u=u_t-a^2u_{xx}=f(x,t), & (x,t)\in Q_T,\\
u(x,0)=\varphi(x), & x\in\mathbb{R}.
\end{cases} \label{eq:heat-cauchy-max1}
\end{align}
的解且满足增长性条件 \eqref{eq:growth}，则
\begin{align*}
\sup_{\overline{Q}_T}|u(x,t)|\leqslant T\sup_{Q_T}|f(x,t)|+\sup_{\mathbb{R}}|\varphi(x)|.
\end{align*}
\end{theorem}
\begin{remark}
由定理我们得到初值问题 \eqref{eq:heat-cauchy-max1}
\begin{align*}
\begin{cases}
\mathcal{L}u=u_t-a^2u_{xx}=f(x,t), & (x,t)\in Q_T,\\
u(x,0)=\varphi(x), & x\in\mathbb{R}.
\end{cases}
\end{align*}
的满足增长性条件 \eqref{eq:growth} 的解的惟一性。然而惟一性在一般情形并不成立。实际上，初值问题
\begin{align*}
\begin{cases}
u_t-u_{xx}=0, & (x,t)\in\mathbb{R}\times(0,+\infty),\\
u(x,0)=0, & x\in\mathbb{R}
\end{cases}
\end{align*}
具有无穷多解。除零解外，每个解在 $|x|\to\infty$ 时都增长得非常快。下面我们引述 Tychonov 的例子来说明上述结论。

记函数
\begin{align*}
\varphi(t)=\begin{cases}
e^{-1/t^2}, & t\neq 0,\\
0, & t=0.
\end{cases}
\end{align*}
定义函数
\begin{align*}
u(x,t)=\begin{cases}
\sum_{n=0}^{\infty}\frac{\mathrm{d}^n}{\mathrm{d}t^n}\varphi(t)\frac{x^{2n}}{(2n)!}, & t\neq 0,\\
0, & t=0.
\end{cases}
\end{align*}
形式上我们有
\begin{align*}
\lim_{t\to 0}u(x,t)=\sum_{n=0}^{\infty}\frac{\mathrm{d}^n}{\mathrm{d}t^n}\varphi(0)\frac{x^{2n}}{(2n)!}=0,
\end{align*}
且
\begin{align*}
\frac{\partial^2 u}{\partial x^2}
&=\sum_{n=1}^{\infty}\frac{\mathrm{d}^n}{\mathrm{d}t^n}\varphi(t)(2n)(2n-1)\frac{x^{2n-2}}{(2n)!} \\
&=\sum_{n=1}^{\infty}\frac{\mathrm{d}^n}{\mathrm{d}t^n}\varphi(t)\frac{x^{2(n-1)}}{(2(n-1))!} \\
&=\sum_{n=0}^{\infty}\frac{\mathrm{d}^{n+1}}{\mathrm{d}t^{n+1}}\varphi(t)\frac{x^{2n}}{(2n)!}=\frac{\partial u}{\partial t}.
\end{align*}
由于篇幅的关系，我们在这里就不给出严格的证明。
\end{remark}
\begin{proof}
实际上我们只需证明不等式
\begin{align*}
|u(0,t)|\leqslant T\sup_{Q_T}|f(x,t)|+\sup_{\mathbb{R}}|\varphi(x)|,\quad t\in[0,T],
\end{align*}
这是因为在平移变换下我们由上述不等式可以得到
\begin{align*}
|u(x,t)|\leqslant T\sup_{Q_T}|f(x,t)|+\sup_{\mathbb{R}}|\varphi(x)|,\quad x\in\mathbb{R},\ t\in[0,T].
\end{align*}
两端取上确界即得要证的结论。

我们分两步来证明：

(1) 首先假设
\begin{align*}
T\leqslant \frac{1}{16a^2A}.
\end{align*}
对于任意 $L>0$，考虑区域 $Q_T^L=(-L,L)\times(0,T]$，并记
\begin{align*}
F=\sup_{Q_T}|f(x,t)|,\qquad \Phi=\sup_{\mathbb{R}}|\varphi(x)|.
\end{align*}
在 $Q_T^L$ 上考虑辅助函数
\begin{align*}
w(x,t)=Ft+\Phi+v_\varepsilon(x,t)\pm u(x,t),
\end{align*}
其中
\begin{align*}
v_\varepsilon(x,t)=\frac{\varepsilon}{(2T-t)^{\frac{1}{2}}}e^{\frac{x^2}{4a^2(2T-t)}},\quad \varepsilon>0.
\end{align*}
注意到
\begin{align*}
\mathcal{L}v_\varepsilon=0,\quad (x,t)\in Q_T^L,
\end{align*}
我们计算得到
\begin{align*}
\begin{cases}
\mathcal{L}w=F\pm f(x,t)\geqslant 0, & (x,t)\in Q_T^L,\\
w|_{t=0}\geqslant \Phi\pm\varphi(x)\geqslant 0, & x\in[-L,L].
\end{cases}
\end{align*}
利用关于 $T$ 的假设和增长性条件 \eqref{eq:growth}，我们得到
\begin{align*}
w|_{x=\pm L}&\geqslant \frac{\varepsilon}{(2T)^{\frac{1}{2}}}e^{\frac{L^2}{8a^2T}}-Me^{AL^2} \\
&\geqslant \frac{\varepsilon}{(2T)^{\frac{1}{2}}}e^{2AL^2}-Me^{AL^2},\quad t\in[0,T].
\end{align*}
对于 $\varepsilon>0$，存在 $L_\varepsilon>0$，使得当 $L\geqslant L_\varepsilon$ 时，
\begin{align*}
w|_{x=\pm L}\geqslant 0,\quad t\in[0,T].
\end{align*}
在 $Q_T^L$ $(L\geqslant L_\varepsilon)$ 上利用极值原理可得 $\min_{\overline{Q}_T^L}w(x,t)\geqslant 0$。对于任意 $t\in[0,T]$，由 $w(0,t)\geqslant 0$ 我们得到
\begin{align*}
|u(0,t)|\leqslant Ft+\Phi+\frac{\varepsilon}{T^{\frac{1}{2}}}.
\end{align*}
令 $\varepsilon\to 0$，则
\begin{align*}
|u(0,t)|\leqslant Ft+\Phi\leqslant FT+\Phi,\quad t\in[0,T].
\end{align*}

(2) 我们将时间区间 $[0,T]$ 分为 $m$ 段：$[0,T_1]$，$[T_1,2T_1]$，$\cdots$，$[(m-1)T_1,mT_1]$，其中 $T=mT_1$ 且
\begin{align*}
T_1\leqslant \frac{1}{16a^2A}.
\end{align*}
先在区间 $[0,T_1]$ 上利用 (1) 的结论从而得到在带型区域 $\mathbb{R}\times[0,T_1]$ 上关于 $u(x,t)$ 的估计，然后在带型区域 $\mathbb{R}\times[T_1,2T_1]$ 上考虑初值问题
\begin{align*}
\begin{cases}
\mathcal{L}u=u_t-a^2u_{xx}=f(x,t), & (x,t)\in\mathbb{R}\times(T_1,2T_1],\\
u|_{t=T_1}=\varphi_1(x)=u(x,T_1), & x\in\mathbb{R},
\end{cases}
\end{align*}
得到在带型区域 $\mathbb{R}\times[T_1,2T_1]$ 上关于 $u(x,t)$ 的估计。利用数学归纳法我们就得到在区域 $\mathbb{R}\times[0,T]$ 上关于 $u(x,t)$ 的估计。

综上我们完成定理的证明.
\end{proof}


\subsection{混合问题的能量模估计}

\begin{lemma}[Gronwall 不等式]
设非负函数 $G(\tau)$ 在 $[0,T]$ 上连续可微，$G(0)=0$，且对 $\tau\in[0,T]$，有微分不等式
\begin{align}
\frac{\mathrm{d}G(\tau)}{\mathrm{d}\tau}\leqslant CG(\tau)+F(\tau), \label{eq:gronwall-diff}
\end{align}
其中 $C>0$，$F(\tau)$ 是 $[0,T]$ 上的非负单调递增函数，那么
\begin{align}
G(\tau)\leqslant C^{-1}(e^{C\tau}-1)F(\tau), \label{eq:gronwall1}
\end{align}
和
\begin{align}
\frac{\mathrm{d}G(\tau)}{\mathrm{d}\tau}\leqslant e^{C\tau}F(\tau). \label{eq:gronwall2}
\end{align}
\end{lemma}
\begin{proof}
在不等式 \eqref{eq:gronwall-diff} 两边同乘 $e^{-C\tau}$，我们得到
\begin{align*}
\frac{\mathrm{d}}{\mathrm{d}\tau}[e^{-C\tau}G(\tau)]\leqslant e^{-C\tau}F(\tau).
\end{align*}
上式两端在 $[0,\tau]$ 上积分得到
\begin{align*}
e^{-C\tau}G(\tau)\leqslant \int_0^\tau e^{-Ct}F(t)\,\mathrm{d}t\leqslant F(\tau)C^{-1}(1-e^{-C\tau}),
\end{align*}
即
\begin{align*}
G(\tau)\leqslant C^{-1}(e^{C\tau}-1)F(\tau).
\end{align*}
将不等式 \eqref{eq:gronwall1} 代入不等式 \eqref{eq:gronwall-diff} 即得不等式 \eqref{eq:gronwall2}.
\end{proof}

\begin{theorem}
设 $u\in C^{1,0}(\overline{Q}_T)\cap C^{2,1}(Q_T)$ 是热方程第一边值问题\eqref{eq:heat-energy}
\begin{align}
\begin{cases}
\mathcal{L}u=u_t-a^2u_{xx}=f, & (x,t)\in Q_T,\\
u|_{t=0}=\varphi(x), & x\in[0,l],\\
u|_{x=0}=u|_{x=l}=0, & t\in[0,T],
\end{cases} \label{eq:heat-energy}
\end{align}
的解，则
\begin{align}
\sup_{0\leqslant t\leqslant T}\int_0^l u^2(x,t)\,\mathrm{d}x+2a^2\int_0^T\int_0^l u_x^2(x,t)\,\mathrm{d}x\,\mathrm{d}t
\leqslant M\left(\int_0^l\varphi^2(x)\,\mathrm{d}x+\int_0^T\int_0^l f^2(x,t)\,\mathrm{d}x\,\mathrm{d}t\right), \label{eq:heat-energy-est}
\end{align}
其中 $M$ 只与 $T$ 有关。
\end{theorem}
\begin{remark}
能量模估计同样适合于其他边值问题。此外，如果在混合问题 \eqref{eq:heat-energy} 的热方程两端同乘 $u_t$，然后在 $Q_t$ 上积分可以得到更进一步的能量模估计
\begin{align*}
a^2\sup_{0\leqslant t\leqslant T}\int_0^l u_x^2\,\mathrm{d}x+\int_0^T\int_0^l u_t^2\,\mathrm{d}x\,\mathrm{d}t
\leqslant 2\left\{a^2\int_0^l(\varphi'(x))^2\,\mathrm{d}x+\int_0^T\int_0^l f^2(x,t)\,\mathrm{d}x\,\mathrm{d}t\right\}.
\end{align*}
\end{remark}
\begin{proof}
在混合问题 \eqref{eq:heat-energy} 的热方程两端同乘 $u$，然后在 $Q_t=(0,l)\times(0,t]$ 上积分，得到
\begin{align*}
\frac{1}{2}\int_0^t\int_0^l [u^2]_t\,\mathrm{d}x\,\mathrm{d}t-a^2\int_0^t\int_0^l u u_{xx}\,\mathrm{d}x\,\mathrm{d}t=\int_0^t\int_0^l uf\,\mathrm{d}x\,\mathrm{d}t.
\end{align*}
对等式左端分部积分并利用边值条件，右端利用 Schwarz 不等式 $2ab\leqslant a^2+b^2$，则
\begin{align*}
\int_0^l u^2(x,t)\,\mathrm{d}x+2a^2\int_0^t\int_0^l u_x^2(x,t)\,\mathrm{d}x\,\mathrm{d}t
\leqslant \int_0^l\varphi^2(x)\,\mathrm{d}x+\int_0^t\int_0^l u^2(x,t)\,\mathrm{d}x\,\mathrm{d}t+\int_0^t\int_0^l f^2(x,t)\,\mathrm{d}x\,\mathrm{d}t.
\end{align*}
记
\begin{align*}
G(t)=\int_0^t\int_0^l u^2(x,t)\,\mathrm{d}x\,\mathrm{d}t,\quad
F(t)=\int_0^l\varphi^2(x)\,\mathrm{d}x+\int_0^t\int_0^l f^2(x,t)\,\mathrm{d}x\,\mathrm{d}t.
\end{align*}
去掉上述不等式左端第二项，我们得到 Gronwall 不等式
\begin{align*}
G'(t)\leqslant G(t)+F(t),
\end{align*}
及
\begin{align*}
G(0)=0.
\end{align*}
利用 Gronwall 不等式的结论，从而得到
\begin{align*}
G(t)\leqslant (e^t-1)F(t).
\end{align*}
将此不等式代入上述不等式得到
\begin{align*}
\int_0^l u^2(x,t)\,\mathrm{d}x+2a^2\int_0^t\int_0^l u_x^2(x,t)\,\mathrm{d}x\,\mathrm{d}t
\leqslant e^t\left(\int_0^l\varphi^2(x)\,\mathrm{d}x+\int_0^t\int_0^l f^2(x,t)\,\mathrm{d}x\,\mathrm{d}t\right).
\end{align*}
对 $t\in(0,T]$ 取上确界得到
\begin{align*}
\sup_{0\leqslant t\leqslant T}\int_0^l u^2(x,t)\,\mathrm{d}x+2a^2\int_0^T\int_0^l u_x^2(x,t)\,\mathrm{d}x\,\mathrm{d}t
\leqslant 2e^T\left[\int_0^l\varphi^2(x)\,\mathrm{d}x+\int_0^T\int_0^l f^2(x,t)\,\mathrm{d}x\,\mathrm{d}t\right].
\end{align*}
这就完成了证明.
\end{proof}

\begin{theorem}
混合问题 \eqref{eq:heat-energy-est1}
\begin{align}
\sup_{0\leqslant t\leqslant T}\int_0^l u^2(x,t)\,\mathrm{d}x+2a^2\int_0^T\int_0^l u_x^2(x,t)\,\mathrm{d}x\,\mathrm{d}t
\leqslant M\left(\int_0^l\varphi^2(x)\,\mathrm{d}x+\int_0^T\int_0^l f^2(x,t)\,\mathrm{d}x\,\mathrm{d}t\right), \label{eq:heat-energy-est1}
\end{align}
的解 $u\in C^{1,0}(\overline{Q}_T)\cap C^{2,1}(Q_T)$ 是惟一的。
\end{theorem}

\begin{theorem}[反向惟一性]
设 $u_1,u_2\in C^1(\overline{Q}_T)\cap C^{2,1}(Q_T)$ 满足
\begin{align}
\begin{cases}
\mathcal{L}u_1=\mathcal{L}u_2, & (x,t)\in(0,l)\times(0,T],\\
u_1|_{t=T}=u_2|_{t=T}, & x\in[0,l],\\
u_1|_{x=0}=u_2|_{x=0},\quad u_1|_{x=l}=u_2|_{x=l}, & t\in[0,T],
\end{cases} \label{eq:backward-unique}
\end{align}
则
\begin{align*}
u_1(x,t)\equiv u_2(x,t),\quad (x,t)\in[0,l]\times[0,T].
\end{align*}
\end{theorem}

\begin{proof}
令 $w=u_1-u_2$，则
\begin{align*}
\begin{cases}
w_t-a^2w_{xx}=0, & (x,t)\in(0,l)\times(0,T],\\
w|_{t=T}=0, & x\in[0,l],\\
w|_{x=0}=0,\quad w|_{x=l}=0, & t\in[0,T]
\end{cases}
\end{align*}
且
\begin{align*}
w_t|_{x=0}=0,\qquad w_t|_{x=l}=0,\quad t\in[0,T].
\end{align*}
记
\begin{align*}
e(t)=\int_0^l w^2(x,t)\,\mathrm{d}x,\quad t\in[0,T],
\end{align*}
则
\begin{align*}
e(T)=0.
\end{align*}
将 $e(t)$ 对 $t$ 求导我们得到
\begin{align*}
e'(t)=2\int_0^l w w_t\,\mathrm{d}x=2a^2\int_0^l w w_{xx}\,\mathrm{d}x=-2a^2\int_0^l w_x^2\,\mathrm{d}x.
\end{align*}
对 $t$ 再求导我们得到
\begin{align*}
e''(t)=-4a^2\int_0^l w_x w_{xt}\,\mathrm{d}x=4a^2\int_0^l w_{xx}w_t\,\mathrm{d}x=4a^4\int_0^l w_{xx}^2\,\mathrm{d}x.
\end{align*}
又由于
\begin{align*}
\int_0^l w_x^2\,\mathrm{d}x=-\int_0^l w w_{xx}\,\mathrm{d}x
\leqslant \left(\int_0^l w^2\,\mathrm{d}x\right)^{\frac{1}{2}}\left(\int_0^l w_{xx}^2\,\mathrm{d}x\right)^{\frac{1}{2}},
\end{align*}
于是
\begin{align*}
(e'(t))^2 &\leqslant 4a^4\left(\int_0^l w^2\,\mathrm{d}x\right)\left(\int_0^l w_{xx}^2\,\mathrm{d}x\right) \\
&=e(t)e''(t).
\end{align*}
对任意 $\varepsilon>0$，记
\begin{align*}
e_\varepsilon(t)=e(t)+\varepsilon.
\end{align*}
显然 $e_\varepsilon(t)\geqslant \varepsilon>0$，于是从上述不等式我们得到
\begin{align*}
[e'_\varepsilon(t)]^2\leqslant e_\varepsilon(t)e''_\varepsilon(t).
\end{align*}
令
\begin{align*}
f(t)=\ln e_\varepsilon(t),\quad 0\leqslant t\leqslant T,
\end{align*}
则
\begin{align*}
f''(t)=\frac{e_\varepsilon(t)e''_\varepsilon(t)-[e'_\varepsilon(t)]^2}{[e_\varepsilon(t)]^2}\geqslant 0,\quad 0\leqslant t\leqslant T.
\end{align*}
因而 $f$ 在区间 $[0,T]$ 上是凸函数。于是对于任意 $\lambda\in(0,1)$，
\begin{align*}
f(\lambda T)\leqslant (1-\lambda)f(0)+\lambda f(T),
\end{align*}
即
\begin{align*}
e_\varepsilon(\lambda T)\leqslant e_\varepsilon(0)^{1-\lambda}e_\varepsilon(T)^\lambda.
\end{align*}
令 $\varepsilon\to 0$，则
\begin{align*}
e(\lambda T)\leqslant e(0)^{1-\lambda}e(T)^\lambda=0.
\end{align*}
由 $\lambda$ 的任意性和 $w(x,t)$ 的连续性我们得到
\begin{align*}
w(x,t)\equiv 0,\quad 0\leqslant t\leqslant T.
\end{align*}
这就完成了证明.
\end{proof}

\subsection{反向问题的不适定性}

在上述内容中我们证明了热方程混合问题和初值问题的适定性，即解的存在惟一性和稳定性。从物理学的观点来看，热传导过程是不可逆的。因此与其逆过程相应的定解问题必定与上述的定解问题有本质的不同。我们将构造一个例子说明这类的定解问题是不适定的。

假设有一根均匀细杆，其侧表面绝热，两端保持零度。已知在 $t=T$ 时刻它的温度分布 $u(x,T)=\varphi(x)$，希望求出 $t=T$ 之前的温度分布 $u(x,t)$。事实上，如果我们希望在 $t=T$ 时刻得到杆的一个理想的温度分布 $\varphi(x)$，就需要知道如何控制杆在 $t=0$ 时刻的初始温度分布来实现。这样的问题称为热方程的\textbf{反向问题}。此时相应于此热传导逆过程的数学问题为
\begin{align}
\begin{cases}
u_t-a^2u_{xx}=0, & (x,t)\in(0,l)\times(0,T),\\
u|_{t=T}=\varphi(x), & x\in[0,l],\\
u|_{x=0}=0,\quad u|_{x=l}=0, & t\in[0,T].
\end{cases} \label{eq:backward-heat}
\end{align}
从之前的结论我们知道，若初值 $u(x,0)$ 为一个连续函数时，解 $u(x,t)$ 在区域 $Q=(0,l)\times(0,T]$ 上是无穷次可微的。事实上，可以证明解 $u(x,t)$ 关于 $x$ 是解析的。因此，如果反向问题 \eqref{eq:backward-heat} 的解存在，则 $\varphi(x)$ 必须是解析函数。换句话说，若 $\varphi(x)$ 仅仅是连续函数，则反向问题 \eqref{eq:backward-heat} 的解肯定不存在。更令人惊讶的是，即使给定的 $\varphi(x)$ 是解析函数且反向问题 \eqref{eq:backward-heat} 的解 $u(x,t)$ 存在，$u(x,t)$ 关于初值 $\varphi(x)$ 也是不稳定的。也就是说，初值 $\varphi(x)$ 的微小误差，也可能引起解 $u(x,t)$ 的巨大改变。而在实际问题中，$\varphi(x)$ 是测量数据，因而测量误差是不可避免的。

下面我们用 Hadamard 的例子来说明此问题。为简单起见我们设 $l=\pi$，$a=1$。令
\begin{align*}
u_k(x,T)=\varphi_k(x)=\frac{1}{k}\sin kx,\quad k=1,2,\cdots,
\end{align*}
容易验证反向问题 \eqref{eq:backward-heat} 的解是
\begin{align*}
u_k(x,t)=\frac{1}{k}e^{k^2(T-t)}\sin kx.
\end{align*}
当 $k\to+\infty$ 时，
\begin{align*}
\max_{[0,\pi]}|\varphi_k(x)|=\frac{1}{k}\to 0.
\end{align*}
但是对于任意 $t<T$，当 $k\to+\infty$ 时，
\begin{align*}
\max_{[0,\pi]}u_k(x,t)=\frac{1}{k}e^{k^2(T-t)}\to+\infty.
\end{align*}
这说明反向问题 \eqref{eq:backward-heat} 在极大模意义下是不适定的。然而，当我们考虑控制问题时，我们会遇到上述反向问题。由于这样的反向问题 \eqref{eq:backward-heat} 具有明确的实际意义，因而我们不得不重新考虑这样的反向问题的适定性问题。为了求解反向问题，我们通常还要加一些适定化条件。这些条件是对解本身所加的一种约束条件，故在这些约束条件下我们仍然可以得到反向问题的适定性。



\end{document}